% Standalone chunk A.3: Minimax subspace-recovery lower bound
% (POST-FIX, Round 2).
% Covers: Thm minimax_recovery (local Grassmannian packing + Fano),
%         Cor id_threshold_conf.
% Source: conf.tex lines ~543-562, 1404-1537 (post-fix).
% Shared preamble for standalone chunks. Do not compile directly; \input{} it.
\documentclass{article}
\usepackage[utf8]{inputenc}
\usepackage[T1]{fontenc}
\usepackage{lmodern}
\usepackage[margin=1in]{geometry}
\usepackage{amsmath,amssymb,amsfonts,amsthm,bm,mathtools}
\usepackage{enumitem}
\usepackage{microtype}
\usepackage{xcolor}
\usepackage{natbib}
\usepackage{booktabs}
\usepackage{hyperref}

\newcommand{\R}{\mathbb{R}}
\newcommand{\N}{\mathbb{N}}
\newcommand{\E}{\mathbb{E}}
\newcommand{\Prob}{\mathbb{P}}
\newcommand{\op}{\mathrm{op}}
\DeclareMathOperator{\sign}{sign}
\newcommand{\cA}{\mathcal{A}}
\newcommand{\cF}{\mathcal{F}}
\newcommand{\cH}{\mathcal{H}}
\newcommand{\cT}{\mathcal{T}}
\newcommand{\cI}{\mathcal{I}}
\newcommand{\cW}{\mathcal{W}}
\newcommand{\cTprobe}{\mathcal{T}_{\mathrm{probe}}}
\newcommand{\cE}{\mathcal{E}}
\newcommand{\norm}[1]{\left\lVert #1 \right\rVert}
\newcommand{\tilO}{\widetilde{\mathcal{O}}}
\newcommand{\DynReg}{\mathrm{DynReg}}
\newcommand{\rank}{\mathrm{rank}}
\newcommand{\range}{\mathrm{range}}
\DeclareMathOperator{\TopEig}{TopEig}
\DeclareMathOperator*{\argmax}{arg\,max}
\DeclareMathOperator{\tr}{tr}

\newtheorem{assumption}{Assumption}
\newtheorem{theorem}{Theorem}
\newtheorem{corollary}{Corollary}
\newtheorem{proposition}{Proposition}
\newtheorem{lemma}{Lemma}
\newtheorem{remark}{Remark}


\title{Chunk A.3: Minimax subspace-recovery lower bound\\(post-fix, Round 2 verification)}
\author{(excerpt for standalone review)}
\date{}

\begin{document}
\maketitle

\textbf{Setup context.}
\begin{itemize}[leftmargin=*,itemsep=1pt,topsep=1pt]
\item Single-segment instance: $B^\star\in\R^{d\times r}$ orthonormal,
$A_1=0$, $w_t\sim\mathcal N(0,\lambda_{\min}I_r)$ i.i.d.,
$\varepsilon_t\sim\mathcal N(0,\sigma_\varepsilon^2)$ i.i.d., probes
$u_t\sim\mathcal N(0,I_d)$ truncated at $\|u\|\le L$.
\item Upper bound (chunk A.2): $\varepsilon_k\le O(\sqrt{\log d/m_k})$.
We show this matches the minimax optimum up to $\sqrt{\log d}$.
\end{itemize}

\textbf{What changed since Round~1 (delta highlights).}
\begin{itemize}[leftmargin=*,itemsep=1pt,topsep=1pt]
\item \emph{Fano constants explicit:} $\alpha = 1/16$ (rather than a
symbolic $\alpha<1/8$), $N\ge 16$ enforced by the packing size
$\log N\ge c_0 r(d-r)$.
\item \emph{Explicit final constant:} $c_1 = \tfrac{1}{32}\sqrt{c_0(d-r)/C_1}$
with $C_1 := 8\lambda_{\min}^2/\sigma_\varepsilon^4$; regime of
validity $m\ge c_0(d-r)/(16 C_1)$ stated.
\item \emph{ID-threshold target risk:} changed to $1/32$ (from $1/4$)
to match the actual constant produced by the Fano lower bound.
\end{itemize}

\section*{Statement}

\begin{theorem}[Minimax subspace-recovery lower bound]
\label{thm:minimax_recovery}
Let $d\ge 2r$ and fix the single-segment model described above with
isotropic Gaussian probes (truncated at $\|u\|\le L$). There exists a
constant $c_1>0$ (depending only on $r,d,L,\sigma_\varepsilon,\lambda_{\min}$,
with $c_1=\Theta(\sqrt{d-r})$) such that for every $m\ge 1$ and every
$\cH_m$-measurable projector estimator $\widehat P$,
\[
\inf_{\widehat P}\ \sup_{\nu\in\mathfrak{M}}\ \E_\nu\!\bigl[\|\widehat P-P^\star\|_{\op}\bigr]
 \ \ge\ \frac{c_1}{\sqrt m}.
\]
The upper bound $O(\sqrt{\log(d/\delta)/m})$ of chunk~A.2 matches this
up to $\sqrt{\log d}$; hence the $T^{2/3}$ probe term in the main
theorem is rate-optimal within $\Pi_{\mathrm{HP}}$.
\end{theorem}

\section*{Proof}

\begin{proof}
The argument is a Fano construction over a \emph{local} Grassmannian
packing with matched upper and lower diameter bounds, so that the KL
bound (upper diameter) and Fano's inequality (lower diameter) are
applied at the same scale.

\emph{Local Grassmannian packing.} For $d\ge 2r$ and any scale
$\Delta\in(0,1/2)$, there exist subspaces $V_1,\dots,V_N\in\mathrm{Gr}(r,d)$
satisfying
\begin{equation}
  \tfrac{1}{2}\Delta \;\le\; \|P_{V_i}-P_{V_j}\|_\op \;\le\; 2\Delta
  \quad\text{for all }i\ne j,
  \qquad \log N \;\ge\; c_0\, r(d-r),
\label{eq:local_packing}
\end{equation}
for an absolute constant $c_0>0$ (Varshamov--Gilbert on a Grassmannian
tangent-space chart around a reference $V_0$; cf.\ Cai--Ma 2013 or
Wainwright HDS Ch.~15). The upper diameter $2\Delta$ is what the KL
computation requires; the lower diameter $\Delta/2$ is what Fano
requires. A global packing with only a separation lower bound would
not suffice.

\emph{Instance construction.} For each $j\in[N]$, let $B_j$ be an
orthonormal basis of $V_j$ and define $\nu_j$: $B^\star=B_j$, $A_1=0$,
$w_t\sim\mathcal N(0,\lambda_{\min}I_r)$ i.i.d.,
$\varepsilon_t\sim\mathcal N(0,\sigma_\varepsilon^2)$ i.i.d.,
$u_t\sim\mathcal N(0,I_d)$ truncated at $\|u\|\le L$ (truncation mass
$\le\delta$ absorbed into constants via Appendix~C chunk). Under $\nu_j$,
the conditional reward distribution given probe $u_t=u$ is zero-mean
Gaussian with variance
$\sigma_j^2(u):=\lambda_{\min} u^\top P_{V_j} u+\sigma_\varepsilon^2\ge\sigma_\varepsilon^2$.

\emph{Pairwise KL.} Let $\Pi_{ij}:=P_{V_i}-P_{V_j}$, so
$\tr\Pi_{ij}=0$. Conditional KL between two zero-mean Gaussians with
variances $\sigma_i^2(u),\sigma_j^2(u)$ is
$\tfrac12(r_u-1-\log r_u)$, $r_u:=\sigma_i^2(u)/\sigma_j^2(u)>0$. Using
the elementary $r-1-\log r\le(r-1)^2/(2\min(r,1))$ and
$\sigma_j^2(u)\ge\sigma_\varepsilon^2$,
\[
\mathrm{KL}(\nu_i^{(y\mid u)}\|\nu_j^{(y\mid u)})
\le\frac{(\sigma_i^2(u)-\sigma_j^2(u))^2}{2\sigma_\varepsilon^4}
=\frac{\lambda_{\min}^2\,(u^\top\Pi_{ij}u)^2}{2\sigma_\varepsilon^4}.
\]
Marginalising over $u\sim\mathcal N(0,I_d)$ (identical probe marginals
under $\nu_i,\nu_j$, so joint KL is
$\E_u[\mathrm{KL}(\nu_i^{(y\mid u)}\|\nu_j^{(y\mid u)})]$), and using
Isserlis $\E[(u^\top M u)^2]=(\tr M)^2+2\|M\|_F^2$ with $\tr\Pi_{ij}=0$,
\[
\E_u[\mathrm{KL}]\le\frac{\lambda_{\min}^2}{2\sigma_\varepsilon^4}\cdot 2\|\Pi_{ij}\|_F^2
=\frac{\lambda_{\min}^2}{\sigma_\varepsilon^4}\|\Pi_{ij}\|_F^2.
\]
$\Pi_{ij}$ is a difference of rank-$r$ projectors, so
$\mathrm{rank}(\Pi_{ij})\le 2r$ and $\|\Pi_{ij}\|_F^2\le 2r\|\Pi_{ij}\|_\op^2\le 8r\Delta^2$ by~\eqref{eq:local_packing}. Hence
\[
\E_u[\mathrm{KL}]\le C_1 r\Delta^2,\qquad
C_1:=\frac{8\lambda_{\min}^2}{\sigma_\varepsilon^4}.
\]
Tensorising over $m$ i.i.d.\ probe rounds,
$\mathrm{KL}(\nu_i^{\otimes m}\|\nu_j^{\otimes m})\le C_1 m r\Delta^2$.

\emph{Fano's inequality.} Standard Fano (Tsybakov Thm.~2.5/Cor.~2.6):
if $\max_{i\ne j}\mathrm{KL}(\nu_i^{\otimes m}\|\nu_j^{\otimes m})\le\alpha\log N$
for $\alpha\in(0,1/8)$ and $N\ge 16$, then
$\inf_{\widehat\psi}\max_j\Pr_{\nu_j}(\widehat\psi\ne j)
\ge\tfrac{\sqrt N}{1+\sqrt N}(1-2\alpha-\sqrt{2\alpha/\log N})\ge 1/4$.
Taking $\alpha=1/16$ suffices for $N\ge 16$; the packing
$\log N\ge c_0 r(d-r)\ge 16$ holds whenever $r(d-r)$ is sufficiently large.
Combining with the lower-diameter bound $\|\Pi_{ij}\|_\op\ge\Delta/2$
via Markov's inequality applied to $\|\widehat P-P_{V_j}\|_\op$,
\[
\inf_{\widehat P}\max_j\E_{\nu_j}[\|\widehat P-P_{V_j}\|_\op]
\ge\frac{\Delta}{2}\cdot\frac{1}{4}=\Delta/8.
\]

\emph{Balancing.} KL constraint $C_1 m r\Delta^2\le\tfrac{1}{16}\log N$,
with $\log N\ge c_0 r(d-r)$, gives $\Delta^2\le c_0(d-r)/(16C_1 m)$.
Setting $\Delta=\min\bigl(1/2,\sqrt{c_0(d-r)/(16C_1 m)}\bigr)$ saturates,
and
\[
\inf_{\widehat P}\sup_{\nu\in\mathfrak M}\E_\nu[\|\widehat P-P^\star\|_\op]
\ge\Delta/8\ge\frac{c_1}{\sqrt m}
\]
with $c_1:=\tfrac{1}{32}\sqrt{c_0(d-r)/C_1}$, valid for
$m\ge c_0(d-r)/(16C_1)$ (regime $m\gg r(d-r)$ where recovery is possible).
\end{proof}

\begin{corollary}[Identifiability sample threshold]
\label{cor:id_threshold_conf}
Any estimator achieving operator-norm subspace error below $1/32$
requires $m\ge c_{\mathrm{id}}(d-r)$ probe rounds.
\end{corollary}
\begin{proof}
Apply Thm.~\ref{thm:minimax_recovery} at target risk $1/32$:
$c_1/\sqrt m\ge 1/32$ needs $m\le 32^2 c_1^2=\Theta(d-r)$ from
$c_1=\Theta(\sqrt{d-r})$. The $\Theta(\sqrt{(d-r)/m})$ op-norm rate is
the standard minimax rate; F-norm gives $\Theta(\sqrt{r(d-r)/m})$,
matching the Grassmannian intrinsic dimension $r(d-r)$.
\end{proof}

\emph{Matching.} Corollary~cor:projector\_conf of chunk~A.2 gives
$\|\widehat P_k-P_k^\star\|_\op\le C_{\mathrm{sub}}\sqrt{\log(2d/\delta)/m_k}$;
Thm.~\ref{thm:minimax_recovery} matches up to the $\sqrt{\log d}$
factor inherent to high-probability matrix Freedman bounds.

\end{document}
